\begin{textsample}{POS Dim 2 – 2012\_c5 – Score -2.00 – 2012\_c5/enade\_2012\_des\_523330154d.txt}  \label{ex:f2_pos_361}
CAPÍTULO 5 DISTRIBUIÇÃO DOS CONCEITOS 5.1 PANORAMA NACIONAL DA DISTRIBUIÇÃO DOS CONCEITOS A Tabela 5.1 apresenta a quantidade e distribuição de cursos de Design participantes do ENADE/2012, por faixa de conceito e Grande Região.
A diferença entre os cursos tabulados neste capítulo e no capítulo 2 são os cursos sem conceito, em princípio, aqueles sem alunos concluintes que participassem da prova ou que tivessem tão somente um aluno concluinte11.
Observando -se os dados da Tabela 5.1, nota -se que, dos 140 cursos participantes, 58 ( 41,4 \% ) classificaram -se com conceito 3, o valor modal.
Este foi também o conceito modal em todas as regiões do Brasil ( entre 40,0 \% no Centro-Oeste e 50,0 \% no Norte ).
O conceito 2 foi o segundo mais frequente em nível nacional ( 27,1 \%, correspondendo a 38 cursos ) e o conceito 4, o terceiro ( 21,4 \%, correspondendo a 30 cursos ).
Houve, ainda, sete cursos ( 5,0 \% ) que receberam conceito 5 e outros três cursos que receberam conceito 1 ( 2,1 \% ).
Apenas quatro dos 140 cursos de Design ficaram sem conceito ( SC ).
A região Norte participou com seis cursos ou 4,3 \% do total nacional.
Nesta região, o conceito 3 foi atribuído a três cursos dos seis participantes, o que equivale a 50,0 \% do total regional, como já comentado.
Nenhum curso ficou sem conceito ( SC ).
Os demais foram avaliados com conceitos 2, 4 e 5 ( um curso em cada, correspondendo a 16,7 \% ).
Nenhum curso foi alocado ao conceito 1.
A região Nordeste participou com 17 cursos ou 12,1 \% do total nacional, destes, sete cursos, 41,2 \% em termos regionais, obtiveram conceito 3, o conceito modal para a região.
Os conceitos 2 e 4 foram atribuídos, igualmente, a cinco cursos ( 29,4 \% ).
Nessa região, tampouco, nenhum dos cursos ficou sem conceito ( SC ) ou recebeu os conceitos 1 ou 5.
Dos 62 cursos participantes da região Sudeste, 25 ( 40,3 \% ) obtiveram conceito 3, o conceito modal.
O conceito 1 foi atribuído a dois cursos ( 3,2 \% ) e o conceito 2, a 21 ( 33,9 \% ).
Nove cursos ( 14,5 \% ) receberam o conceito 4 e dois outros ( 3,2 \% ) receberam o conceito 5.
Nesta região três cursos ficaram sem conceito ( SC ).
A região Sul também contou com cursos distribuídos em todas as faixas de conceitos.
A predominância do conceito 3 foi de 42,0 \%, correspondentes a 21 dos 50 cursos participantes na região Sul.
O conceito 2 foi atribuído a dez cursos ( 20,0 \% ) e o conceito 4, a 14 cursos ( 28,0 \% ).
O conceito 1 foi recebido por um único curso ( 2,0 \% ) e o conceito 5, por quatro deles ( 8,0 \% ).
Nenhum dos cursos da região Sul ficou sem conceito ( SC ).
A região Centro-Oeste participou com apenas cinco cursos.
Como já comentado, dois destes receberam conceito 3 ( 40,0 \% ), o conceito modal.
Nesta região um curso ficou sem conceito ( SC ).
Os demais cursos foram avaliados com conceito 2 e 4 ( um curso cada, 20,0 \% ).
Nenhum curso foi alocado nos conceitos 1 ou 5. 5.2 CONCEITOS POR CATEGORIA ADMINISTRATIVA E POR GRANDE REGIÃO A Tabela 5.2 apresenta a distribuição dos cursos participantes do ENADE/2012 de Design, por Categoria Administrativa, de acordo com os conceitos por eles alcançados, segundo as Grandes Regiões brasileiras.
Dos 140 cursos participantes, 35 ( 25,0 \% ) eram ministrados em instituições públicas e 105 ( 75,0 \% ), em privadas.
De acordo com as informações da Tabela 5.2, em termos nacionais, entre instituições públicas e privadas : das sete IES que receberam o conceito 5, cinco eram públicas e dois privadas.
Dos 35 cursos participantes de IES públicas, o conceito 4 foi o valor modal, atribuído a 14 cursos ( 40,0 \% ).
Entre os demais cursos de instituições públicas participantes, quatro obtiveram conceito 2 ( 11,4 \% da categoria ), 12 receberam conceito 3 ( 34,4 \% ) e como já comentado, os cinco restantes foram alocados ao conceito 5 ( 14,3 \% ).
Nesta categoria, nenhum dos cursos ficou sem conceito ( SC ) ou recebeu o conceito 1.
Na rede privada, o conceito modal foi 3, com 46 cursos dos 105 da categoria.
Nesta categoria administrativa, quatro dos cursos ficaram sem conceito ( SC ).
Entre os demais cursos participantes, três receberam conceito 1 e 34, conceito 2.
O conceito 4 foi atribuído a 16 cursos e o conceito 5, como já comentado, a dois cursos. 11 Estes cursos com somente um concluinte foram considerados no capítulo 2.
Os conceitos 2 e 4 foram atribuídos, igualmente, a cinco cursos ( 29,4 \% ).
Nessa região, tampouco, nenhum dos cursos ficou sem conceito ( SC ) ou recebeu os conceitos 1 ou 5.
Dos 62 cursos participantes da região Sudeste, 25 ( 40,3 \% ) obtiveram conceito 3, o conceito modal.
O conceito 1 foi atribuído a dois cursos ( 3,2 \% ) e o conceito 2, a 21 ( 33,9 \% ).
Nove cursos ( 14,5 \% ) receberam o conceito 4 e dois outros ( 3,2 \% ) receberam o conceito 5.
Nesta região três cursos ficaram sem conceito ( SC ).
A região Sul também contou com cursos distribuídos em todas as faixas de conceitos.
A predominância do conceito 3 foi de 42,0 \%, correspondentes a 21 dos 50 cursos participantes na região Sul.
O conceito 2 foi atribuído a dez cursos ( 20,0 \% ) e o conceito 4, a 14 cursos ( 28,0 \% ).
O conceito 1 foi recebido por um único curso ( 2,0 \% ) e o conceito 5, por quatro deles ( 8,0 \% ).
Nenhum dos cursos da região Sul ficou sem conceito ( SC ).
A região Centro-Oeste participou com apenas cinco cursos.
Como já comentado, dois destes receberam conceito 3 ( 40,0 \% ), o conceito modal.
Nesta região um curso ficou sem conceito ( SC ).
Os demais cursos foram avaliados com conceito 2 e 4 ( um curso cada, 20,0 \% ).
Nenhum curso foi alocado nos conceitos 1 ou 5. 5.2 CONCEITOS POR CATEGORIA ADMINISTRATIVA E POR GRANDE REGIÃO A Tabela 5.2 apresenta a distribuição dos cursos participantes do ENADE/2012 de Design, por Categoria Administrativa, de acordo com os conceitos por eles alcançados, segundo as Grandes Regiões brasileiras.
Dos 140 cursos participantes, 35 ( 25,0 \% ) eram ministrados em instituições públicas e 105 ( 75,0 \% ), em privadas.
De acordo com as informações da Tabela 5.2, em termos nacionais, entre instituições públicas e privadas : das sete IES que receberam o conceito 5, cinco eram públicas e dois privadas.
Dos 35 cursos participantes de IES públicas, o conceito 4 foi o valor modal, atribuído a 14 cursos ( 40,0 \% ).
Entre os demais cursos de instituições públicas participantes, quatro obtiveram conceito 2 ( 11,4 \% da categoria ), 12 receberam conceito 3 ( 34,4 \% ) e como já comentado, os cinco restantes foram alocados ao conceito 5 ( 14,3 \% ).
Nesta categoria, nenhum dos cursos ficou sem conceito ( SC ) ou recebeu o conceito 1.
Na rede privada, o conceito modal foi 3, com 46 cursos dos 105 da categoria.
Nesta categoria administrativa, quatro dos cursos ficaram sem conceito ( SC ).
Entre os demais cursos participantes, três receberam conceito 1 e 34, conceito 2.
O conceito 4 foi atribuído a 16 cursos e o conceito 5, como já comentado, a dois cursos.
Na análise por região, observa -se que, na região Norte, as instituições privadas participaram com cinco cursos ( 83,3 \% do total regional ), dos quais nenhum ficou sem conceito ( SC ).
O conceito modal para as instituições privadas na região foi 3, com três cursos, correspondendo a 60,0 \%.
Os outros dois cursos receberam conceitos 2 e 4 ( um Na região Nordeste, a rede privada concentrou oito dos 17 cursos participantes, equivalentes a 47,1 \% do total da região, a menor proporção entre todas as regiões.
Nesta Categoria Administrativa na região, predominaram os cursos com conceito 2, quatro cursos correspondendo a exata metade dos cursos oferecidos por IES privadas no Nordeste.
Os demais foram avaliados com conceitos 3 ( três cursos correspondendo a 37,5 \% ) e 4 ( um curso correspondendo a 12,5 \% ).
Nenhum curso recebeu o conceito 1, nem tampouco o conceito 5.
As instituições públicas dessa região participaram com nove cursos ( 52,9 \% ), dos quais quatro cada obtiveram conceitos 3 e 4, os valores modais.
O outro curso foi avaliado com o conceito 2.
Nenhum curso ficou sem conceito ( SC ) nesta combinação de região e categoria.
Na região Sudeste, a proporção de cursos da rede privada, 83,9 \%, foi mais elevada do que nas demais regiões brasileiras, correspondendo a 52 dos 62 cursos participantes.
Nesta categoria, na região Sudeste, o conceito modal foi 3 ( 22 cursos ).
Nesta combinação de Categoria Administrativa e Grande Região, três cursos não receberam conceito ( SC ) e nenhum foi alocado ao conceito 5.
Os demais foram avaliados com conceito 1 ( dois cursos ), conceito 2 ( 19 cursos ) e conceito 4 ( seis cursos ).
Entre os dez cursos em instituições públicas na região Sudeste, as categorias modais foram as que receberam conceitos 3 e 4 ( três cursos cada ).
Os demais foram avaliados com os conceitos 2 e 5 ( dois cursos em cada ).
Na região Sudeste, nesta categoria, nenhum curso ficou sem conceito ( SC ).
As instituições privadas concentraram 37 dos 50 cursos participantes da região Sul, 74,0 \% do total regional.
Desses, 17 ficaram com conceito 3, o conceito modal.
Os demais foram avaliados com conceito 1 ( um curso ), conceito 2 ( nove cursos ), conceito 4 ( oito cursos ) e conceito 5 ( dois cursos ).
Nesta combinação de Categoria Administrativa e Grande Região, nenhum curso ficou sem receber conceito ( SC ).
As instituições públicas na região Sul participaram com 13 cursos ( 26,0 \% ), dos quais nenhum ficou sem conceito.
O conceito modal foi 4, com seis cursos.
Além disso, o conceito 2 foi recebido por um curso, o conceito 3 por quatro cursos e conceito 5 por dois cursos.
Nenhum curso foi alocado ao conceito 1.
Na região Centro-Oeste, três dos cinco cursos participantes eram de instituições privadas ( 60,0 \% em termos regionais ).
Destes, um curso ficou sem conceito ( SC ) e os outros dois se distribuíram igualmente nos conceitos 2 e 3.
Dos dois cursos de instituições públicas, um foi avaliado no conceito 3 e o outro no conceito 4.
Nesta região, nenhum curso de IES pública ficou sem conceito ( SC ). 5.3 CONCEITOS POR ORGANIZAÇÃO ACADÊMICA E POR GRANDE REGIÃO Na Tabela 5.3 encontra -se a distribuição dos conceitos atribuídos aos cursos participantes do ENADE/2012 na Área de Design, por Organização Acadêmica, segundo as Grandes Regiões brasileiras.
Dos 140 cursos participantes, 70 eram oferecidos em Universidades, 23 em Centros Universitários e os demais 47 em Faculdades.
Esta distribuição corresponde a, respectivamente, 50,0 \%, 16,4 \% e 33,6 \% dos cursos.
De acordo com os dados apresentados, de todos os sete cursos avaliados com conceito 5, seis eram vinculados a Universidades.
Apenas três cursos em Universidades ficaram sem conceito ( SC ).
Esse tipo de Organização Acadêmica teve o conceito 3 como modal, com 26 cursos ( 37,1 \% ).
Os demais cursos de Universidades avaliados receberam os conceitos 1 ( um curso ), 2 ( 12 cursos ), 4 ( 22 cursos ) e conceito 5 ( seis cursos, como já mencionado ).
Entre os cursos em Centros Universitários, o conceito modal também foi 3, com 14 cursos ( 60,9 \% ).
Neste tipo de organização acadêmica apenas um curso ficou sem conceito ( SC ), enquanto nenhum curso recebeu o conceito 4.
Os outros cursos neste tipo de Organização Acadêmica receberam os conceitos 1 ( um curso ), 2 ( seis cursos ) e 5 ( também um curso ).
Nas Faculdades, nenhum dos 47 cursos ficou sem conceito ( SC ) e 20 ( 42,6 \% ) receberam o conceito modal, 2.
Dos demais cursos neste tipo de Organização Acadêmica, um recebeu conceito 1, 18 conceito 3 e oito cursos, o conceito 4.
Nenhum curso recebeu o conceito 5.
Considerando -se separadamente as regiões brasileiras, verifica -se que, na região Norte, as Universidades concentraram somente um dos seis cursos participantes, este Os Centros Universitários da região Norte não foram representados por nenhum curso.
As Faculdades participaram com cinco cursos na região Norte.
Nenhum ficou sem conceito ( SC ) ou recebeu os conceitos 1 ou 5.
Os demais receberam conceito 2 ( um curso ), 3 ( três cursos ) e conceito 4 ( também um curso ).
Na região Nordeste, as Universidades participaram com 11 dos 17 cursos na Área de Design da região.
Nenhum dos cursos em Universidades no Nordeste ficou sem conceito ( SC ) ou recebeu conceitos 1 ou 5.
Os conceitos modais foram 3 e 4, com quatro cursos em cada.
Os demais receberam conceito 2 ( três cursos ).
Os Centros Universitários contaram com somente um curso participante na região Nordeste, este avaliado com conceito 3.
As Faculdades foram representadas por cinco cursos na região Nordeste.
Os conceitos modais foram 2 e 3, com dois cursos cada, enquanto o outro curso foi avaliado com conceito 4.
Na região Sudeste, as Universidades concentraram 26 dos 62 cursos da região.
Entre os cursos em Universidades na região, o conceito modal foi 3 com 11 cursos e apenas dois cursos ficaram sem conceito ( SC ).
Os demais cursos receberam os conceitos 1 ( um curso ), 2 ( seis cursos ), 4 ( quatro cursos ) e 5 ( dois cursos ).
Os Centros Universitários participaram com 13 cursos na região Sudeste, dos quais oito obtiveram conceito modal, 3, e apenas um ficou sem conceito ( SC ).
Os demais receberam os conceitos 1 ( um curso ) e 2 ( três cursos ).
Nenhum dos cursos nesta combinação de organização acadêmica e região recebeu conceitos 4 ou 5.
As Faculdades foram representadas por 23 cursos na região Sudeste, que foram distribuídos nos conceitos 2 ( 12 cursos, conceito modal ), 3 ( seis cursos ) e 4 ( cinco cursos ).
Nenhum curso ficou sem conceito ( SC ).
Dos 50 cursos da região Sul, 28 eram de Universidades, para os quais o conceito modal foi 4, com 13 cursos.
Nesse tipo de organização, nenhum dos cursos ficou sem conceito ( SC ) e os demais receberam os conceitos 2 ( três cursos ), 3 ( nove cursos ) e 5 ( três cursos ).
Os Centros Universitários da região Sul tiveram cinco dos nove cursos participantes no conceito modal, 3.
Os outros quatro cursos receberam conceitos 2 ( três cursos ) e 5 ( um curso ).
Foram 13 os cursos vinculados a Faculdades na região Sul, e nenhum destes ficou sem conceito ( SC ).
Um desses cursos recebeu conceito 1, quatro receberam conceito 2, sete receberam conceito 3, o conceito modal, e outro curso recebeu conceito 4.
Na região Centro-Oeste, quatro dos cinco cursos eram de Universidades.
Nesse tipo de organização, apenas um curso ficou sem conceito ( SC ) e o conceito modal foi 3, com dois cursos.
O outro curso obteve o conceito 4.
Os Centros Universitários da região Centro-Oeste não contaram com nenhum curso.
As Faculdades da região Centro-Oeste participaram com somente um curso, que foi avaliado com conceito 2.

% matched lemmas: 
\end{textsample}
